% BackdoorLens paper. All numbers/tables/findings are \input from ../results/<run>/report (generated by scripts/make_report.py).
% Build: scripts/build_paper.sh <results-dir>.  Uses IEEEtran when installed (e.g. Overleaf), otherwise a two-column article fallback.
\IfFileExists{IEEEtran.cls}{\documentclass[conference]{IEEEtran}\newcommand{\hasieee}{1}}{\documentclass[10pt,twocolumn]{article}\newcommand{\hasieee}{0}}
\usepackage[utf8]{inputenc}
\usepackage[T1]{fontenc}
\usepackage{booktabs,amsmath,amssymb,graphicx,xcolor,url,array,microtype}
\ifnum\hasieee=0
  \usepackage[margin=0.75in,top=0.8in]{geometry}\usepackage{mathptmx}\usepackage{cite}
  \newenvironment{IEEEkeywords}{\par\noindent\textbf{Keywords---}}{\par}
  \newcommand{\IEEEauthorblockN}[1]{#1}\newcommand{\IEEEauthorblockA}[1]{\\\small #1}
\fi
\usepackage[colorlinks=true,linkcolor=blue!50!black,citecolor=blue!50!black,urlcolor=blue!50!black]{hyperref}
\input{/tmp/toy1/report/macros.tex}
\newcommand{\bl}{\textsc{BackdoorLens}}
\newcommand{\JSD}{\mathrm{JSD}}
\def\/tmp/toy1/report{/tmp/toy1/report}
\begin{document}
\title{Transformation-Aware Behavioral Attestation for Derived Open-Weight LLM Artifacts: Calibrated Counterexample Search and Deployment-Scoped Integrity Enforcement}
\author{\IEEEauthorblockN{Ayush Munot}\IEEEauthorblockA{MCA, Savitribai Phule Pune University, Maharashtra, India}}
\maketitle
\ifnum\syntheticrun=1
\begin{center}\fbox{\parbox{0.92\columnwidth}{\color{red}\textbf{SYNTHETIC VALIDATION BUILD.} The numbers below come from a toy / random-weight model used only to test the pipeline. They are \emph{not} research evidence. Re-run \texttt{scripts/run\_experiments.py} on a real model and rebuild before submission.}}\end{center}
\fi
\begin{abstract}
Vetting an open-weight language model at its source does not establish that the artifact finally served---after quantization, format conversion, adapter merging, tokenizer or chat-template packaging, generation-configuration changes and runtime optimization---behaves like the vetted reference. Signatures, file scanners, perplexity and aggregate benchmarks verify identity, file safety and average quality, not conditional behavior. We present \bl{}, a white-box verifier that compares a vetted reference bundle with the exact derived deployment bundle. It combines (i) a deterministic bundle differential over tokenizer, sandbox-rendered chat template and configuration; (ii) paired next-token divergence measurement; (iii) a benign-divergence null fitted per transformation family and tested with a derivation-aware statistic; (iv) a budgeted counterexample search compared with random and fixed-benchmark probes at identical cost; and (v) a signed, replayable attestation that is invalidated automatically when any in-scope component changes. We evaluate on controlled harmless tampering (template injection, generation-config edits, planted conditional behavior, weight deltas) with honest-derivation controls. On our \nHonestAudits{} held-out honest audits the false-alarm rate was \fprGuided\% (95\% CI upper bound \fprGuidedHi\%), and \invTP{} of \invN{} in-scope modifications invalidated the certificate. A passing attestation is not a proof of safety: it is bounded evidence under a declared scope, input distribution and budget.
\end{abstract}
\begin{IEEEkeywords}
LLM supply chain, behavioral attestation, backdoors, quantization, model integrity, calibration
\end{IEEEkeywords}

\section{Introduction}
Open-weight models rarely reach production in the form in which they were vetted. A deployment chain may include quantization, conversion between runtime formats, adapter merging, tokenizer and chat-template packaging, generation-configuration changes and compiler or kernel optimization. Recent work shows that each stage can matter for security: quantization can activate behavior that is dormant in full precision~\cite{qtb,agentq}, inference optimization can act as an activation condition~\cite{opttrig}, and a chat template---an executable Jinja program---can implant conditional behavior without changing any weight~\cite{chattmpl}.

Three trust questions are routinely conflated: \emph{provenance} (is this the file I expected?), \emph{file safety} (does it contain malicious code?) and \emph{behavioral faithfulness} (does it behave like the reference I vetted?). Model signing~\cite{sigstore} addresses the first; hub scanners and template scanners address parts of the second; \bl{} targets the third and treats the others as complementary inputs.

\textbf{Contributions.} We keep the claims narrow and individually testable: (C1) a \emph{scope tuple} that binds evidence to weights, adapters, tokenizer, template, configuration, generation policy, quantization and runtime, with explicit unknown/excluded states; (C2) a \emph{bundle differential} that renders untrusted templates in a sandbox on a structured conversation corpus; (C3) a \emph{benign-divergence null} decomposed into level and shape, with a derivation-aware test for the search maximum that avoids pseudo-replication (Sec.~\ref{sec:null}); (C4) a \emph{matched-budget comparison} of benchmark, random and guided counterexample search; (C5) a \emph{signed, replayable attestation} with deterministic enforcement and automatic invalidation. We do not claim universal backdoor detection, a proof of safety, or that guided search outperforms simpler probes: Sec.~\ref{sec:results} reports what was measured, including negative outcomes.

\section{Background and Related Work}
\textbf{Backdoor detection.} ToxScreen~\cite{toxscreen} benchmarks about 800 backdoored models under white-box access without a reference and reports that no method surfaces every backdoor. Backdoor Probes~\cite{bprobes} scan with forward passes; Watch the Weights~\cite{wtw} reads fine-tuning deltas from weight differences when a base model exists; SpecGuard~\cite{specguard} uses speculative-decoding signals at inference time. These detect implanted behavior; \bl{} instead asks whether a \emph{derived} artifact stays faithful to a vetted reference and uses such detectors as optional reference-vetting.
\textbf{Derivation-conditioned attacks and defenses.} Quantization-triggered~\cite{qtb,agentq}, optimization-triggered~\cite{opttrig} and template-based~\cite{chattmpl} attacks motivate scoping the evidence to the served bundle. Reported defenses focus on remediation, whereas we address verification.
\textbf{Equality and attestation.} Model equality testing~\cite{met} and IRIS~\cite{iris} verify what an API serves under black-box access. The closest work to ours, adversarial probes with privacy-preserving proofs~\cite{zkattest}, searches for probes that amplify drift between an approved and a modified deployment; we differ in the white-box bundle scope, the transformation-matched null with a derivation-aware search test, matched-budget baselines and lifecycle invalidation. We have read only its abstract and treat the overlap as a risk to be resolved before submission.
\textbf{Adaptive auditing and monitoring.} Anytime-valid adaptive auditing~\cite{adaptive} provides sequential inference for audits; we use its statistical discipline (predeclared tests, separated splits) but make no anytime-validity claim. NIST AI 800-4~\cite{nist} names post-deployment monitoring as a gap; our lifecycle module is one concrete mechanism.

\section{Threat Model and Scope}
Let $R=(W,A,T,C,G,Q,X)$ denote the reference weights, adapters, tokenizer and template, configuration, generation policy, quantization state and runtime metadata, and $D$ the corresponding derived bundle. We estimate whether divergence between $P_R(y\mid x)$ and $P_D(y\mid x)$ over a \emph{declared input distribution} is abnormal for a \emph{declared transformation family} under a finite audit budget; we never claim $P_R=P_D$ for all $x$.
\textbf{Actors.} A malicious publisher; a compromised conversion or quantization pipeline; a maintainer who edits the template, tokenizer or generation configuration; a supply-chain actor who swaps a package after certification; and an adaptive attacker who knows the probe policy.
\textbf{Trust.} The reference is conditionally trusted (it must be vetted independently); the audit engine is the trusted computing base; the derived bundle is untrusted until attested; an LLM, if used to summarize, is non-authoritative and never sets a verdict.
\textbf{Out of scope.} Proprietary black-box models; automated purification; universal safety claims. If the reference itself is poisoned, a faithful derivation preserves the poison.

\section{Method}
\subsection{Scope tuple and bundle differential}
Each component is hashed and labeled \emph{known}, \emph{unknown} (never guessed) or \emph{excluded}. A certificate is valid only while every in-scope component matches. The bundle differential compares tokenizer files and special tokens, configuration and generation fields, and the \emph{rendered} output of the reference and derived chat templates on twelve structured conversations (system, multi-turn, tool-like, empty, unicode, special-token text). Templates are untrusted executable code and are rendered only in an immutable sandboxed Jinja environment.

\subsection{Paired measurement}
For each prompt both bundles use their own template and tokenizer. We compute $\JSD(P_R,P_D)$ over next-token distributions in bits (primary), plus KL, top-$k$ overlap and rank correlation, margin and NLL shifts, and generated-text structural flips and edit distance as the output-only baseline.

\subsection{Benign null and decision statistic}\label{sec:null}
\textbf{Level and shape.} Honest derivations of one family differ in overall drift (e.g.\ INT4 with different group alignment). Treating that variation as evidence of tampering inflates false alarms. For a derived artifact we therefore test (i) its \emph{level} $\ell$, the median per-prompt JSD, against the spread of levels of $k$ honest derivations with a $t$ prediction interval on $\log\ell$ ($k-1$ degrees of freedom), and (ii) its \emph{shape}, per-prompt JSD and search maximum divided by $\ell$.
\textbf{Search null.} A percentile of per-prompt JSD is not a valid threshold for a \emph{searched maximum}: searching inflates it. We therefore run the identical search procedure at the identical budget on honest derivations over many replicates and store the maxima of $\mathrm{best}/\ell$.
\textbf{Avoiding pseudo-replication.} Replicates on one honest derivation are not independent evidence about a \emph{new} derivation. We model $\log m_{dj}=\mu+u_d+\epsilon_{dj}$ and use the prediction variance $s_b^2+s_w^2$ with $\mathrm{SE}^2=(s_b^2+s_w^2/R)/k$ and $k-1$ degrees of freedom. Sec.~\ref{sec:results} reports the false-alarm rate of the naive pooled-replicate rule next to this one. The log-normal assumption is checked empirically by leave-one-derivation-out audits on the calibration split, not assumed.
\textbf{Effect-size guard.} A strong verdict additionally requires the counterexample's ratio to exceed a margin times the largest honest ratio, so that statistical tails alone cannot drive a block.

\subsection{Budgeted counterexample search}
Three strategies receive an identical budget of unique paired evaluations (one reference and one derived forward pass each): \emph{benchmark} (curated seeds and rule-based variants), \emph{random} (seed plus random mutations) and \emph{guided} (exploration followed by tournament selection of parents by $J(x)=\lambda_1 A(x)+\lambda_2 F(x)-\lambda_3 B(x)$, where $A$ is $\log_2(1+\JSD/\ell)$, $F$ an argmax flip and $B$ an out-of-distribution penalty). Mutations are controlled lexical, structural, role-format and position edits; validity is enforced by rules, not self-certified by the target model.

\subsection{Confirmation and verdicts}
The top candidates are re-executed (replay), and tested on four rule-based semantic variants; a counterexample is \emph{persistent} if at least three variants exceed the honest per-prompt null and replay matches. Verdicts are produced by a deterministic policy: INVALIDATED, INCONCLUSIVE (unknown transformation, no matching null, too few honest derivations, insufficient budget), BEHAVIORALLY\_DIVERGENT, SUSPICIOUS, ATTESTED\_WITH\_DRIFT, ATTESTED. A passing verdict is never labeled ``safe''.

\subsection{Attestation and lifecycle}
The attestation records the scope tuple, hashes, probe-suite identifier, seeds, budget, null hash, metrics, counterexamples with replay data, policy version, environment and expiry, and is signed with Ed25519. Signing authenticates the record; the evidence is the measurement. A gate maps the certificate state to CI exit codes and a lightweight canary fingerprint detects unexplained drift.

\section{Experimental Protocol}
\textbf{Run.} Backend: \texttt{\backend}; probe suite \texttt{\suiteId} (128 prompts in six families: general, safety, template, format, position, instruction-pressure); budget \budget{} paired evaluations; calibration, test and tampered artifacts are disjoint.
\textbf{Honest derivations.} Simulated symmetric round-to-nearest group quantization at 4 and 8 bits with varied group alignment. This reproduces rounding error but is not a bitsandbytes, GPTQ or GGUF kernel; real-toolchain derivations are future work.
\textbf{Controlled tampering (harmless).} Hidden instruction in the chat template; generation-configuration edit; planted conditional behavior producing a canned marker (an in-lexicon trigger and an out-of-lexicon trigger); weight-noise canaries of increasing magnitude.
\textbf{Baselines.} Hash and signature (passes by construction), serialization scan, template static scan, perplexity-style NLL check, output-only text drift, global thresholds, pooled-replicate search rule, and benchmark/random search at equal budget.
\textbf{Statistics.} Calibration and thresholds are fitted only on the calibration split; test artifacts are audited once per seed; CIs are exact binomial and optimistic where audits of one artifact are correlated. All numbers are generated from \texttt{results.json}.

\section{Results}\label{sec:results}
\subsection{Computed findings}
\input{/tmp/toy1/report/findings.tex}
\subsection{Tables and figures}
\input{/tmp/toy1/report/tables/t1_conventional.tex}
\input{/tmp/toy1/report/tables/t2_fpr.tex}
\begin{figure}[t]\centering\includegraphics[width=\columnwidth]{/tmp/toy1/report/figures/fig_detection.pdf}\caption{Detection rate by strategy at identical paired-case budget (95\% CI).}\label{fig:det}\end{figure}
\input{/tmp/toy1/report/tables/t3_search.tex}
\input{/tmp/toy1/report/tables/t4_calibration.tex}
\begin{figure}[t]\centering\includegraphics[width=0.8\columnwidth]{/tmp/toy1/report/figures/fig_levels.pdf}\caption{Honest benign drift level differs by transformation family, so a single global threshold is not appropriate.}\label{fig:lev}\end{figure}
\input{/tmp/toy1/report/tables/t5_output_vs_logit.tex}
\input{/tmp/toy1/report/tables/t6_invalidation.tex}

\section{Discussion}
\textbf{What the evidence supports.} Only the statements in Sec.~\ref{sec:results} are claims of this paper. Negative outcomes (artifacts no strategy detected, any strategy ties) bound the contribution and are reported as such.
\textbf{Where attestation helps and where it does not.} The bundle differential is deterministic and exact for template, tokenizer and configuration changes. Behavioral search is distribution- and budget-scoped; rare triggers outside the exercised distribution can evade it, and an adaptive attacker who knows the probe policy can place behavior outside it (mitigations: randomized seeds, held-out probes, an evasion study).
\textbf{Reproducible rebuild.} If a recipe is public and deterministic, rebuilding and comparing hashes is stronger than any behavioral test; \bl{} is for derived artifacts whose rebuild is unavailable or non-deterministic.

\section{Limitations and Ethics}
(1) Faithfulness is not innocence: a poisoned reference is preserved. (2) Finite probes cover only the declared distribution, family and budget. (3) The null requires several honest derivations of the same family; otherwise the verdict is INCONCLUSIVE. (4) Quantization is simulated in the reported run; real toolchains may differ. (5) Planted behaviors are harmless canned markers; no harmful payload is produced, and adaptive-attack results should be reported at the level of outcomes. (6) Small models only; scaling behavior is not measured here. (7) The log-normal prediction interval is an assumption checked by leave-one-derivation-out audits.

\section{Conclusion}
\bl{} measures whether a deployed derived bundle remains behaviorally faithful to a vetted reference under a declared scope, calibrates what honest derivations look like, preserves replayable evidence and enforces a deterministic admission decision. The measured evidence in Sec.~\ref{sec:results} is the contribution; it does not establish that a passing model is safe.

\begin{thebibliography}{99}
\bibitem{qtb} Dardini et al., ``Quantization-triggered backdoors and the validation--deployment gap,'' arXiv:2608.27512, 2026.
\bibitem{agentq} Liu and Yan, ``AgentQ,'' arXiv:2609.14060, 2026.
\bibitem{opttrig} Wang et al., ``Optimization-triggered backdoor attacks,'' arXiv:2605.20641, 2026.
\bibitem{chattmpl} Fogel et al., ``Inference-time backdoors via chat templates,'' arXiv:2602.04653, 2026.
\bibitem{sigstore} Sigstore, ``Model transparency v1.0,'' 2025. [Online]. Available: https://blog.sigstore.dev/model-transparency-v1.0/
\bibitem{toxscreen} Hughes, Xing, Francel, Kim, Draganov, ``ToxScreen: Detecting whether an LLM has been poisoned,'' arXiv:2607.26849, 2026.
\bibitem{bprobes} Bullwinkel et al., ``The trigger in the haystack,'' arXiv:2602.03085, 2026.
\bibitem{wtw} Zhong and Raghunathan, ``Watch the weights,'' ICLR 2026, arXiv:2508.00161.
\bibitem{specguard} Wen, Salem, Paverd, Russinovich, Li, ``SpecGuard: Inference-time backdoor detection for free,'' arXiv:2609.11799, 2026.
\bibitem{met} ``Model equality testing,'' arXiv:2410.20247.
\bibitem{iris} ``IRIS,'' arXiv:2607.20860, 2026.
\bibitem{zkattest} ``Not to break, but to attest: Adversarial probes for privacy-preserving LLM verification,'' arXiv:2608.27954, 2026.
\bibitem{adaptive} Zhou, Vossler, Sivaraman, Mai, Feng, ``Adaptive auditing of AI systems with anytime-valid guarantees,'' arXiv:2605.07002, 2026.
\bibitem{nist} NIST, ``Challenges to the monitoring of deployed AI systems,'' NIST AI 800-4, Mar. 2026.
\end{thebibliography}
\end{document}
